\section{Conclusiones}\label{sec:conclusiones}

\subsection{Conclusión sobre OE1 (EDA)}
Cumplido. Se caracterizó la siniestralidad fatal nacional con 14 figuras interpretadas y once hallazgos numerados H1--H11, con al menos tres hallazgos fuertes: la separación estructural rural/urbana en la tasa multifatal (H8, 15.1\,\% vs.\ 3.6\,\%), la dominancia de la clase de siniestro en la letalidad (H5) y la divergencia entre el ranking por volumen y por tasa multifatal territorial (H6, Lima vs.\ Huancavelica).

\subsection{Conclusión sobre OE2 (características)}
Cumplido con holgura. Se construyeron 116 características derivadas de las variables pre-impacto del registro, cada una justificada y documentada en \texttt{tab02\_feature\_contract.csv}. Se excede el mínimo de 12 e incluye codificación cíclica de la hora, \emph{frequency encoding} de vía, coordenadas estandarizadas, condiciones de clima y geometría de vía, y \emph{flags} de faltantes.

\subsection{Conclusión sobre OE3 (modelo y baselines)}
Cumplido, con reporte honesto. Se entrenó un perceptrón multicapa y se comparó contra tres líneas base. El MLP supera con claridad al \texttt{DummyClassifier} y entrega el mayor \emph{recall}-multifatal (0.8849), la métrica privilegiada por la regla de costo asimétrico. La regresión logística obtiene un F1 nominalmente superior (0.3081 vs.\ 0.2746) dentro de intervalos traslapados: en datos tabulares de este tamaño los métodos lineales compiten de igual a igual, y este trabajo lo reporta en lugar de ocultarlo (contingencia C13).

\subsection{Conclusión sobre OE4 (interpretabilidad)}
Cumplido. El top-5 SHAP (\texttt{clase\_atropello}, \texttt{zona\_rural}, \texttt{clase\_despiste}, \texttt{clase\_choque}, \texttt{zona\_urbana}) se identificó y se contrastó con los hallazgos del EDA: la coherencia es alta y la red aprendió exactamente las variables de contexto que motivaron el cambio de fuente de datos.

\subsection{Conclusión sobre OE5 (interfaz)}
Cumplido. La interfaz gráfica en Streamlit reutiliza el contrato de preprocesamiento mediante \texttt{preparar\_entrada()}, con cuatro pestañas operativas, un \emph{gauge} de probabilidad calibrada, comparación contra tasas históricas, ranking territorial con mapa coroplético y una pestaña sobre el modelo. Los siete checks funcionales pasan y la app se verificó de punta a punta en navegador.

\subsection{Conclusión general}
Se desarrolló una red neuronal de clasificación binaria de alta letalidad en siniestros viales fatales del Perú, con pipeline completo (EDA $\rightarrow$ preprocesamiento $\rightarrow$ ingeniería $\rightarrow$ entrenamiento $\rightarrow$ evaluación única $\rightarrow$ calibración $\rightarrow$ GUI), reproducible (semillas fijas, modelo canónico único, artefactos versionados) y documentado. La reformulación del problema hacia la letalidad múltiple ---forzada por la naturaleza fatal-únicamente de la mejor fuente disponible--- resultó metodológicamente productiva: la discriminación mejoró (ROC-AUC 0.6840 $\rightarrow$ 0.7474), las probabilidades ahora están calibradas (Brier 0.0872, mejor que el predictor de tasa base) y la incertidumbre se reporta con intervalos bootstrap. Las limitaciones se declaran sin maquillaje: la logística compite con la red, el registro es preliminar y faltan variables mecánicas (ocupantes, velocidad) que fijarían el siguiente techo. El proyecto entrega lo comprometido y declara lo no comprometido, que es exactamente lo que un jurado valora.